% !TeX root = ../../tesis.tex

% =============================================================================
%  B.9 — GUÍA DE REPLICABILIDAD
% =============================================================================
\section{Guía de Replicabilidad}
\label{anx:vft-replicabilidad}

Esta sección condensa lo necesario para levantar VFTModel en un entorno local
desde cero. El código fuente del repositorio está disponible en
\url{https://github.com/galigaribaldi/VFTModel}.

\subsection*{Stack tecnológico}

\begin{table}[H]
    \centering
    \small
    \caption[Stack tecnológico de VFTModel]{Stack tecnológico requerido para
    ejecutar VFTModel.
    \fuentefigura{Elaboración propia (VFTModel, 2026).}}
    \label{tab:vft-stack}
    \begin{tabular}{|l|l|}
        \hline
        \textbf{Componente}   & \textbf{Versión} \\ \hline
        Lenguaje              & Python~3.12 \\ \hline
        API framework         & FastAPI~0.110.0 + Uvicorn~0.29.0 \\ \hline
        Validación            & Pydantic~2.6.4 \\ \hline
        Cliente HTTP          & httpx~0.27.0 \\ \hline
        Análisis de grafos    & NetworkX~3.2.1 \\ \hline
        Análisis espacial     & GeoPandas~0.14.3 + Shapely~2.0.3 \\ \hline
        Morfología urbana     & Momepy~0.7.0 \\ \hline
        Sistema operativo dep.& GDAL (\texttt{libgdal-dev} / \texttt{gdal-bin}) \\ \hline
        Contenedores          & Docker (imagen \texttt{python:3.12-slim-bookworm}) \\ \hline
        Documentación \api{}  & Swagger UI en \texttt{/docs} \\ \hline
        Tests                 & pytest (39 tests de integración) \\ \hline
    \end{tabular}
\end{table}

\subsection*{Levantar el entorno local (sin Docker)}

\begin{lstlisting}[language=bash,
    caption={Secuencia completa para levantar VFTModel en desarrollo local.},
    label=lst:vft-setup]
# 1. Clonar el repositorio
git clone https://github.com/galigaribaldi/VFTModel.git
cd VFTModel

# 2. Crear y activar entorno virtual
python3 -m venv venv
source venv/bin/activate   # macOS / Linux

# 3. Instalar dependencias (requiere GDAL instalado previamente)
make install

# 4. Configurar la URL de Apimetro
cp .env.example .env.local
# Editar .env.local: APIMETRO_URL=http://localhost:8080/movilidad
\end{lstlisting}

\subsection*{Configuración multi-entorno}

VFTModel distingue dos entornos de conexión mediante archivos \texttt{.env}:

\begin{lstlisting}[language=bash,
    caption={Selección de entorno mediante variable \texttt{ENV\_FILE}.},
    label=lst:vft-envfile]
# Desarrollo local (Apimetro corriendo en localhost:8080)
make run                         # lee .env.local por defecto

# Desarrollo contra servidor remoto (Apimetro en apimetro.dev)
make run-dev                     # lee .env.dev
# o bien:
ENV_FILE=.env.dev make run-dev
\end{lstlisting}

\subsection*{Secuencia de cálculo: del grafo a los indicadores}

Los endpoints de indicadores exigen que el grafo esté construido y en caché. La
secuencia correcta es:

\begin{lstlisting}[language=bash,
    caption={Secuencia mínima para calcular indicadores tras levantar el servidor.},
    label=lst:vft-sequence]
# Terminal 1: levantar el servidor
make run

# Terminal 2: construir el grafo y calentarlo en cache
curl "http://localhost:8000/api/v1/network/build-auto?mode=REALISTIC_INTEGRATION"
# Respuesta esperada: {"status":"success","nodos":11115,"aristas":45679}

# Una vez que build-auto confirmo exito:
curl "http://localhost:8000/api/v1/network/spatial-coverage"
curl "http://localhost:8000/api/v1/network/topological/capillary-strength"
curl "http://localhost:8000/api/v1/network/topological/detour-factor"
curl "http://localhost:8000/api/v1/network/topological/average-travel-time"
curl "http://localhost:8000/api/v1/network/topological/betweenness-centrality"
\end{lstlisting}

La documentación interactiva de todos los endpoints con sus parámetros está
disponible en \texttt{http://localhost:8000/docs} (Swagger UI generado
automáticamente por FastAPI).

\subsection*{Levantar con Docker}

\begin{lstlisting}[language=bash,
    caption={Secuencia para construir y correr VFTModel en Docker.},
    label=lst:vft-docker]
# Construir la imagen (una sola vez)
make docker-build

# Docker apuntando al servidor DEV remoto (apimetro.dev)
make docker-run PORT=8000

# Docker apuntando a Apimetro local (Linux — requiere host-gateway)
make docker-run-local PORT=8000
\end{lstlisting}

\subsection*{Correr la suite de pruebas}

La suite cuenta con 39 pruebas de integración que verifican la cadena completa:
validación Pydantic, construcción del grafo y cálculo de cada indicador.
Requiere el servidor activo y acceso a Apimetro:

\begin{lstlisting}[language=bash,
    caption={Ejecución de la suite pytest de VFTModel.},
    label=lst:vft-tests]
make run-dev   # Terminal 1: servidor apuntando a apimetro.dev
make test      # Terminal 2: pytest tests/ -v
\end{lstlisting}
